\begin{definition}
    \textbf{Prodotto scalare}

    Sia $\mathbb{X}$ uno spazio vettoriale su $\mathbb{R}$, allora una qualsisasi operazione
    \[\langle \cdot, \cdot \rangle: \mathbb{X}\times\mathbb{X}\rightarrow\mathbb{R}\]
    viene detta prodotto scalare o interno se gode delle seguenti proprietà.
    \begin{property}
        \textbf{Positività}
        \[
            \begin{cases}
                \forall\overline{x}\in\mathbb{X},\;
                \langle\overline{x},\overline{x}\rangle \geq 0\\
                \forall\overline{x}\in\mathbb{X},\;
                \langle\overline{x},\overline{x}\rangle = 0 \Leftrightarrow
                \overline{x}=\overline{0}
            \end{cases}.
        \]
    \end{property}
    \begin{property}
        \textbf{Simmetria}
        \[
            \forall\overline{x},\overline{y}\in\mathbb{X},\;
            \langle\overline{x},\overline{y}\rangle=\langle\overline{y},\overline{x}\rangle.
        \]
    \end{property}
    \begin{property}
        \textbf{Additività}
        \[
            \forall\overline{x},\overline{y},\overline{z}\in\mathbb{X},\;
            \langle\overline{x}+\overline{y},\overline{z}\rangle=
            \langle\overline{x},\overline{z}\rangle +
            \langle\overline{y},\overline{z}\rangle.
        \]
    \end{property}
    \begin{property}
        \textbf{Omogeneità}
        \[
            \forall\overline{x},\overline{y}\in\mathbb{X},\, \alpha\in\mathbb{R},\;
            \alpha\langle\overline{x},\overline{y}\rangle = \langle\alpha\overline{x},\overline{y}
            \rangle = \langle\overline{x},\alpha\overline{y}\rangle.
        \]
    \end{property}
    \begin{example}
        \[
            \text{\underline{Spazio vettoriale:}} \quad \mathbb{X }= \left\{f:\left[a, b\right]\rightarrow\mathbb{R} \mid f \text{ è integrabile e continua}\right\}
        \]\[
            \text{\underline{Prodotto scalare:}} \quad \langle f,g\rangle = \int_{a}^{b}f\cdot g \;dx,
        \]
    \end{example}
\end{definition}

\begin{definition}
    \textbf{Norma}

    Sia $\mathbb{X}$ uno spazio vettoriale su $\mathbb{R}$, allora una qualsisasi funzione
    \[\| \cdot \|: \mathbb{X}\rightarrow\mathbb{R}\] su
    viene detta norma se gode delle seguenti proprietà.
    \begin{property}
        \textbf{Positività}
        \[
            \begin{cases}
                \forall\overline{x}\in\mathbb{X},\;
                \|\overline{x}\| \geq 0\\
                \forall\overline{x}\in\mathbb{X},\;
                \|\overline{x}\| = 0 \Leftrightarrow
                \overline{x}=\overline{0}
            \end{cases}.
        \]
    \end{property}
    \begin{property}
        \textbf{Omogeneità assoluta}
        \[
            \forall\overline{x}\in\mathbb{X},\, \alpha\in\mathbb{R},\;
            \left|\alpha\right|\cdot\|\overline{x}\| = \|\alpha \overline{x}\|.
        \]
    \end{property}
    \begin{property}
        \textbf{Disuguaglianza triangolare}
        \[
            \forall\overline{x},\overline{y}\in\mathbb{X},\;
            \|\overline{x}+\overline{y}\| \leq
            \|\overline{x}\| + \|\overline{y}\|.
        \]
    \end{property}
\end{definition}

\begin{proposition}\label{pro:induzione-norma}
    La definizione di un prodotto scalare in uno spazio
    vettoriale induce una norma deinfita come segue.
    \[
        \| \overline{x} \| := \sqrt{\langle \overline{x}, \overline{x} \rangle}.
    \]
\end{proposition}
Per dimostrare quest'ulitma proposizione occorre per comodità dimostrare prima la seguente.

\begin{theorem}\label{pro:cauchy-schwarz}
    \textbf{Disuguaglianza di Cauchy-Schwarz}

    Sia $\mathbb{X}$ un qualsiasi spazio vettoriale su $\mathbb{R}$, sul quale è definito il prodotto
    vettoriale $\langle \cdot, \cdot \rangle$, allora
    \[
        \forall \overline{x}, \overline{y}\in\mathbb{X},
        \left|\langle \overline{x}, \overline{y} \rangle\right| \leq
        \|\overline{x}\|\cdot\|\overline{y}\|.
    \]
\end{theorem}

\begin{demonstration}
    \textit{di} \teoref{pro:cauchy-schwarz}
    
    [dimostrazione]
\end{demonstration}

\begin{demonstration}
    \textit{di} \proref{pro:induzione-norma}

    [dimostrazione]
\end{demonstration}